\chapter{Programs and Lexical Structure}
\label{ch:lexical}

\section{Source files}

An \algolang\ program is a single UTF-8 text file with the extension \file{.algo}. It consists of top-level declarations in any order:

\begin{center}
\small
\begin{tabularx}{\linewidth}{@{}L{2.6cm}L{4.2cm}Y@{}}
\toprule
\thead{Declaration} & \thead{Form} & \thead{Purpose}\\
\goldrule
\kw{model} & \code{model name = "id"} & Names a model identifier\\
\kw{tool} & \code{tool web.search} & Declares a capability agents may be granted\\
\kw{agent} & \code{agent N(p: t) -> t \{ … \}} & Declares a unit of work\\
\kw{flow} & \code{flow main(p: t) \{ … \}} & Orchestrates agents\\
\bottomrule
\end{tabularx}
\end{center}

Every program must contain exactly one \code{flow main}, which is where execution starts (error \errcode{E111} otherwise).

\section{Comments, whitespace and semicolons}

\begin{itemize}
  \item A comment starts with \code{//} and runs to the end of the line.
  \item Whitespace, including newlines, separates tokens and is otherwise ignored. Indentation has no meaning.
  \item Semicolons are allowed and ignored, so \code{a = A(x); b = B(a)} on one line is the same as two lines.
\end{itemize}

\section{Identifiers}

Identifiers match \code{[a-zA-Z\_][a-zA-Z0-9\_]*}. By convention agents use \code{PascalCase} (\code{Researcher}), and variables, parameters and models use \code{lower\_case} (\code{facts}, \code{smart}). Names of agents and flows share one namespace; defining the same name twice is error \errcode{E107}.

Tool names are \emph{dotted} identifiers with at least two parts, such as \code{web.search} or \code{file.read}.

\section{Keywords}

\begin{center}
\small
\begin{tabularx}{\linewidth}{@{}L{3.3cm}Y@{}}
\toprule
\thead{Group} & \thead{Keywords}\\
\goldrule
Declarations & \kw{model} \kw{tool} \kw{agent} \kw{flow}\\
Agent clauses & \kw{role} \kw{goal} \kw{uses} \kw{model} \kw{policy}\\
Policy rules & \kw{budget} \kw{tokens} \kw{retry} \kw{timeout} \kw{require} \kw{human} \kw{when} \kw{confidence} \kw{forbid}\\
Flow statements & \kw{par} \kw{for} \kw{in} \kw{if} \kw{else} \kw{emit}\\
Literals & \kw{true} \kw{false}\\
Types & \kw{text} \kw{score} \kw{list} \kw{json} \kw{bool}\\
\bottomrule
\end{tabularx}
\end{center}

\section{Literals}

\begin{itemize}
  \item \textbf{Strings} are enclosed in double quotes. Backslash escapes such as \code{\textbackslash n} and \code{\textbackslash"} are supported.
  \item \textbf{Numbers} are written as \code{6000} or \code{0.75}. Negative and exponent forms are not part of version 0.1.
  \item \textbf{Booleans} are \kw{true} and \kw{false}.
\end{itemize}

\section{Placeholders in goals}

Inside a \kw{goal} string, \code{\{name\}} is a placeholder for the parameter \code{name}. When the agent runs, the placeholder is replaced by the argument's value (a list is joined with \code{"; "}). The compiler checks that every placeholder names a real parameter (\errcode{E109}).

\begin{algo}
agent Translator(passage: text, language: text) -> text {
    goal "Translate into {language}, keeping technical terms: {passage}"
}
\end{algo}

\section{Operators}

\begin{center}
\small
\begin{tabularx}{\linewidth}{@{}L{3.3cm}Y@{}}
\toprule
\thead{Operator} & \thead{Use}\\
\goldrule
\code{->} & Separates an agent's parameters from its return type\\
\code{=} & Assignment in flows; default values in parameters; keyword arguments\\
\code{<} \code{<=} \code{>} \code{>=} \code{==} \code{!=} & Comparisons in \kw{if} conditions\\
\code{[ , ]} & Binds several results of a \kw{par} block\\
\bottomrule
\end{tabularx}
\end{center}

\algolang\ has no arithmetic or string operators. Computation is done by agents; flows only route values between them.
